% ─────────────────────────────────────────────────────────────────────────────
\subsection{Σύνθεση: Πέντε Κλειδιά Ευρήματα}
\label{sec:fi_synthesis}
% ─────────────────────────────────────────────────────────────────────────────

Η ανάλυση σημαντικότητας χαρακτηριστικών αναδεικνύει πέντε διακριτά μοτίβα
που επεξηγούν γιατί κάθε στρατηγική συμπεριφέρεται όπως συμπεριφέρεται.

\begin{enumerate}

\item \textbf{TF: αυτο-αναδρομική κυριαρχία.}
Στη Teacher-Forcing στρατηγική, το \textit{Price\_lag1} συγκεντρώνει
${\sim}88\%$ της σημαντικότητας gain για το LightGBM τιμής. Το μοντέλο
έχει μάθει ότι η τιμή της προηγούμενης ώρας είναι ο ισχυρότερος προβλεπτής
της επόμενης. Αυτό είναι στατιστικά αποδοτικό, αλλά δημιουργεί εύθραυστα
μοντέλα σε distributional shift, όπου η αυτο-αναδρομική δυναμική αλλάζει
ρυθμό (βλ. κατάρρευση SVR σε Q1 2026).

\item \textbf{Rec: κατανεμημένη σημαντικότητα, ανάδυση εξωγενών.}
Η Recursive στρατηγική χρησιμοποιεί τις ίδιες ιδιότητες χωρίς πρόσβαση σε πραγματικές ενδοημερήσιες τιμές.
Το \textit{Price\_lag1} καταρρέει σε ${\sim}9\%$, με αυξημένη βαρύτητα στα
\textit{Price\_lag24} (ημερήσια εποχικότητα) και \textit{Price\_lag168} (εβδομαδιαία).
Για το XGB-SS-Dense, η Scheduled Sampling αναγκάζει το μοντέλο να αξιοποιεί
τα dense intraday lags (4–23) ως επιπλέον πηγή πληροφορίας.

\item \textbf{MIMO: στροφή στα εξωγενή.}
Χωρίς πρόσβαση σε ενδιάμεσες πραγματικές τιμές, το MIMO-LGBM αναδεικνύει στις
πρώτες θέσεις τις αιολικές υστερήσεις (\textit{wind\_gen} lag24/48/168) και τον
κυλιόμενο μέσο. Για το RF-MIMO two-stage,
το \textit{Price\_roll24} (κυλιόμενος μέσος) αναλαμβάνει ρόλο anchor.
Για το XGB-MIMO Dense, τα dense lags \textit{Price\_lag4}–\textit{Price\_lag23}
αποκτούν συλλογικά ${\sim}53\%$ σημαντικότητας, αύξηση που εξηγεί τη
βελτίωση $-7{,}9\%$ MAE έναντι του XGB-MIMO standard.

\item \textbf{Dense lags paradox: MIMO vs Rec.}
Τα dense intraday lags βελτιώνουν το MIMO ($-22\%$ για MLP, $-7{,}9\%$ XGB,
$-6{,}5\%$ LGBM) αλλά επιβαρύνουν το Rec ($+2{,}8\%$ LGBM). Αιτία:
στο MIMO δεν υπάρχει recursion, οπότε τα lags είναι πάντα καθαρές
ιστορικές τιμές. Στο Rec, αντίθετα, τα dense lags γεμίζουν από
ήδη-υπολογισμένες προβλέψεις: θόρυβος που επεκτείνεται στα επόμενα βήματα.
Το Scheduled Sampling αντιμετωπίζει μερικώς αυτό (LGBM-SS-Dense: $-1{,}6\%$
έναντι LGBM-SS).

\item \textbf{Φορτίο: χαμηλότερη συγκέντρωση, υψηλότερη ρολή εποχικότητας.}
Για τη στρατηγική TF φορτίου, το \textit{Load\_lag1} συγκεντρώνει ${\sim}70\%$
(έναντι $88\%$ για τιμή). Η σημαντικότητα κατανέμεται πλουσιότερα: ημερολογιακά
χαρακτηριστικά (ώρα, ημέρα εβδομάδας), \textit{Load\_lag24} (κυρίως στο
RandomForest, ${\sim}17\%$) και \textit{Load\_roll24} αποκτούν σημαντικές θέσεις. Η ομαλότερη δυναμική
του φορτίου επιτρέπει στο Recursive XGB-SS να ανταγωνίζεται στατιστικά την
TF ($p = 0{,}769$, DM test Q1 2026).

\end{enumerate}
